% 원제: 3. 설계 요구사항과 판정 기준
\chapter{설계 요구사항과 판정 기준}

\section{요구사항}

요구사항은 기존 보고에서 확정한 것을 기본으로 하고, 개선 구조와 이번 검증에서 도출한 하위 요구를 더했다(표 \ref{t:req}). 400±8 V 정밀 조정은 외부 부하의 요구가 아니라 고정이득 Stage C와 B\textsubscript{420} 부품 호환에서 나오는 파생 요구였다\cite{b9}, \cite{b10}. R2에서는 외부 요구인 버스 창 조건(부하 하한 이상, 부품 등급 상한 이하)을 직접 기준으로 쓴다.

{\footnotesize
\begin{longtable}{L{\dimexpr 0.2193\linewidth-2\tabcolsep\relax}L{\dimexpr 0.5351\linewidth-2\tabcolsep\relax}L{\dimexpr 0.2456\linewidth-2\tabcolsep\relax}}
\caption{설계 요구사항과 판정 기준}\label{t:req}\\
\toprule
\tabhead \textbf{항목} & \textbf{값} & \textbf{근거·상태} \\
\midrule
\endfirsthead
\toprule
\tabhead \textbf{항목} & \textbf{값} & \textbf{근거·상태} \\
\midrule
\endhead
\bottomrule
\endlastfoot
배터리 운전범위 & 사양 팩 640–920 V. 팩 등급 검토는 셀 180–240s(3.0–4.2 V) & E4, E18 \\
400 V 버스 창 & 부하 하한(260 V) 이상, 부품 등급 상한(B\textsubscript{420} 420 V, B\textsubscript{470} 470 V) 이하. 창 조정 모듈은 V\textsubscript{W}(410·460 V) 기준 & 규격\cite{b9} \\
창 조정 전압 품질 & Tier-1: V\textsubscript{W} ±8 V·리플 ≤ 0.1 V / Tier-2: 상한 이하 & E13·E19 \\
연속전력·예산 & 25 kW(62.5 A), SALS 예산 23.75 kW & 부하 등급·계약 \\
전류 판정 기준 & 평균모델 Tier-1 90 A·Tier-2 93.75 A / 스위칭·회로 모델(리플 포함 순시) 98·100 A & 물리정격 100 A \\
제어 지연 & 모델 내 5 µs 영차유지 + 5 µs 전달지연, 모델 밖 예산 +3.0 µs & E6·E7 \\
창 조정 DAB 정격 & 정적 최대 링크전력 × 1.8(R2-W460 4.0 kW, R2-W410 5.0 kW) & 도출(E19) \\
센서·ADC & 16비트, 버스 350–450 V 차동, 잡음 표준편차 전압 5 mV·전류 20 mA & 도출(E12) \\
게이트 시간 분해능 & H-브리지 PWM·DAB 위상 에지 ≤ 0.625 ns(고분해능 PWM) & 도출(E12, TB04B) \\
충전전류 증가 램프 & ≥ 100 µs(375 A 기준 ≤ 3.75 A/µs) & 도출(E16, TB07) \\
고장 강등 & 경로·필터 ≤ 100 A, 직렬·링크 ≤ 120 V, 바이패스 ≤ 329 A, 버스 ≥ 260 V, 400 V 부하 연결 중 버스 ≤ 상한, 강등 상태 도달 & E17·TB09 \\
강등 중 DAB 첨두 & 링크 쪽 순시 첨두 ≤ 1.5 × 정격점 첨두(R2-W460 68.7 A, R2-W410 90.0 A) & 도출(E20) \\
보호 하드웨어 & 링크 제동 저항 3 Ω(고장당 ≤ 64 J), 바이패스 게이트 선형 영역 구동, 400 V 부하 분리 경로 & 도출(E17) \\
400 V급 급속충전 / 48·12 V 전원 & 피크 150 kW·375 A 제한 / 48 V 3 kW + 12 V 1.5 kW 동시(Stage C 5 kW), 48 V 회생 −3 kW & 부하 등급 \\
조정형 Stage C & 입력 270–504 V, 48 V ±2 \%(계약 과도 ≤ 4 kW/ms)·36–52 V(계단 견고성), 출력 정격 105 A(과도 130 A, 2 ms), 5 kW 효율 ≥ 96 \% & 규격\cite{b10}, E22 \\
Stage C 데드타임 & 1차 100 ns·2차 40 ns(두 브리지 구동 때 90 ns). 출력용량 ±20 \%에서 통과하는 창의 가운데이며 창 폭은 20–30 ns & 도출(E22, TB14·TB15) \\
Stage A 영전류 스위칭 & 턴오프 전류/피크 ≤ 5 \%(출력 7.5 A 미만은 턴오프 전류 ≤ 1 A), 위상 첨두 ≤ 76.4 A(정격), 모듈 1개 고장 시 ≤ 114.6 A(1.5배) & E1, E24, E26 \\
Stage A 위상 셰딩 & 멈춤 15 A(100 µs), 복귀 22 A 또는 순시 36 A, 자연 턴오프 뒤 멈춤 & 도출(E24) \\
충전기 비상 정지 & 정지 신호 뒤 30 ms 안에 출력 ≤ 5 A, 차량 접촉기는 5 A 이하에서 개방, 커넥터 ≤ 600 V & 규격\cite{b16} \\
V\textsubscript{A} 클램프 & 468 V 켬·463 V 끔, 1 Ω, 스위치 지연 ≤ 1 µs, 펄스 7 kJ·첨두 220 kW(충전 1회 고장당), 비교기 정확도 약 ±0.1 \% & 도출(E26, TB20) \\
Stage A 위상당 고정 손실 & 현재 모델(15.6 W)에서 늘어나는 몫 < 14.8 W(합 약 30 W 미만). 넘으면 모터 승압 구조 대비 보조 손실 이점이 없어진다. 공진 인덕터는 공심(철심 손실 규칙이면 위상당 +18.8 W) & 도출(E25) \\
\end{longtable}
}
\srcnote{주: 근거 번호는 부록 C의 검증 근거를 따른다.}

\section{가설별 수락 기준}

가설은 표 \ref{t:accept}의 기준으로 판정한다. 기준은 요구사항에서 도출해 시험 전에 고정했다. H5는 기존 보고의 기준 그대로 판정했고, 그 결과(기각)에 따라 적용 범위를 묻는 판정을 더했다. 현재 판정 상태는 15.1절에 정리한다.

{\footnotesize
\begin{longtable}{L{\dimexpr 0.2105\linewidth-2\tabcolsep\relax}L{\dimexpr 0.3509\linewidth-2\tabcolsep\relax}L{\dimexpr 0.4386\linewidth-2\tabcolsep\relax}}
\caption{가설별 판정 지표와 수락 기준}\label{t:accept}\\
\toprule
\tabhead \textbf{가설(질문)} & \textbf{판정 지표} & \textbf{수락 기준} \\
\midrule
\endfirsthead
\toprule
\tabhead \textbf{가설(질문)} & \textbf{판정 지표} & \textbf{수락 기준} \\
\midrule
\endhead
\bottomrule
\endlastfoot
H1 정격 분리(Q1) & Stage B 링크전력, 충전전류 경로 & 링크 ≤ 8 kW(정상·Tier-1·Tier-2, 공차 포함), 충전전류 Stage B 비경유 \\
H2 강건 제어(Q2) & Tier-1·Tier-2, 파라미터 ±10 \% 공차, 지연 예산 & Tier-1: 전압 품질 + 전류 기준. Tier-2: 물리정격 + 전류 기준 + 버스 창. 공차와 추가 지연 3.0 µs에서 Tier-1 유지 \\
H3 충전 중 보조버스(Q1) & 150 kW(375 A 제한) 충전 + 25 kW 보조부하, 충전기 접속·해제 & 동시 운전에서 Tier-1, 모드 전환에서 Tier-2 \\
H4 부하 승인(Q2) & 부하 계약, 메시지 지연 1–5 ms, 손실 0–20 \% & 예산 초과 0, 컴프레서 승인 지연 ≤ 25 ms, 미공급이 정적 예약보다 작을 것 \\
H5 기준 구조 대비 이점(Q3) & 기준 구조 ①·②·④와 변형 ③′ 대비 동일 경계 비교 & 손실·정격 VA·반도체 수 중 2개 이상에서 ④와 ③′보다 우수, 서비스 동등 이상 \\
Q3 적용 범위(추가) & 셀 180–240s × 부품 B\textsubscript{420}·B\textsubscript{470}의 성립성과 손실 & 성립 구조 가운데 최소 손실인 팩·부품 조합이 있을 것 \\
Q3 산업 방식 대비(추가) & 모터·인버터 승압 충전(⑤), 800 V급 보조부품 설계(⑥) 대비 손실·반도체 수 & R2가 앞서는 조건과 뒤지는 조건을 같은 손실 규칙으로 구분할 것 \\
\end{longtable}
}
\srcnote{주: ① 전전력 승압기 + 800→400 V 강압기, ② Stage A + 전전력 조정기, ③′ 창 조정 변형(버스 ≤ 410 V), ④ 비조정(버스 = V\textsubscript{A,} R2 기본형), ⑤ 모터·인버터 승압 충전 + 800→400 V 강압기, ⑥ 모터·인버터 승압 충전 + 800 V급 보조부품.}
